\mode*
\section{Related work}

\Textcite{IsomottonenCochez2014} document Git learning challenges from
course experience and recommend visual, pointer-based teaching of
branches; repository-data studies of student Git use add that students
commit rarely and treat the system as file storage
\autocite{Cochez2013DVCSUsage}, and a large-scale Stack Overflow study
shows the command-level difficulty is intrinsic to Git rather than peculiar
to novices \autocite{Yang2022GitCommands}.
\Textcite{DeRossoJackson2016Gitless} locate part of the difficulty in Git's
concept design and demonstrate a redesign (Gitless).
On the teaching side, the GitKit line catalogues misconceptions in basic
GitHub workflows across institutions and packages an authentic FOSS project
for learning them
\autocite{Postner2026GitMisconceptions,Braught2024GitKit};
the Shell Tutor tutors Git through the shell \autocite{Winder2024ShellTutor};
and visualisation tools make the hidden state discernible, from commit-graph
diagrams to a split-view collaborative simulator
\autocite{Bucher2026GitTakesTwo}.
Pedagogical experience points the same way our analysis does: students
report being overwhelmed by Git's many commands and options, and teaching
is advised to convey the underlying Git model rather than command recipes
\autocite{Wagner2025RethinkGit}---independent corroboration of the
notional-machine argument.
Our contribution is the variation-theory analysis on top of this literature:
from documented misconceptions to critical aspects to teachable patterns of
variation, as in the companion papers
\autocite{VTProgMisconceptions,VTDebug,VTTerminal}.
